% =====================================================================
%  1.1 Potenser — teori og oppgaver   (MØNSTERSEKSJON)
% =====================================================================

\del{1}{Potensreglene}

\begin{ordliste}
grunntall & Tallet som ganges med seg selv. I $2^5$ er grunntallet 2. & \mittsprak\\
eksponent & Hvor mange ganger grunntallet er faktor. I $2^5$ er eksponenten 5. & \mittsprak\\
potens & Et uttrykk på formen $a^n$. & \mittsprak\\
\end{ordliste}

\begin{mustkunne}
\begin{minipage}[t]{0.47\linewidth}
\textbf{Definisjon.} $a^n$ betyr $n$ like faktorer:
\[ a^n = \underbrace{a \cdot a \cdot \ldots \cdot a}_{n \text{ faktorer}} \]
\textbf{Hvorfor regelen virker:}
\[ a^3 \cdot a^2 = (a\cdot a\cdot a)\cdot(a\cdot a) = a^5 \]
Vi teller faktorer. Derfor legger vi sammen eksponentene.
\end{minipage}\hfill
\begin{minipage}[t]{0.49\linewidth}
\textbf{Potensreglene}\par
\renewcommand{\arraystretch}{1.5}
\begin{tabular}[t]{@{}l l@{}}
$a^m \cdot a^n = a^{m+n}$ & Produkt\\
$\frac{a^m}{a^n} = a^{m-n}$ & Kvotient\\
$(a^m)^n = a^{m \cdot n}$ & Potens av potens\\
$(a \cdot b)^n = a^n \cdot b^n$ & Potens av produkt\\
$\left(\frac{a}{b}\right)^n = \frac{a^n}{b^n}$ & Potens av brøk\\
\end{tabular}
\end{minipage}
\end{mustkunne}

\begin{eksempel}[— forenkle steg for steg]
Forenkle $\dfrac{(2a^3)^2 \cdot a}{4a^5}$.

\begin{stegtabell}
$= \dfrac{4a^6 \cdot a}{4a^5}$ & Potens av produkt: $2^2 = 4$ og $(a^3)^2 = a^6$.\\
$= \dfrac{4a^7}{4a^5}$ & Produkt: $a^6 \cdot a^1 = a^7$.\\
$= a^2$ & Kvotient: $a^{7-5} = a^2$. Og $\frac{4}{4} = 1$.\\
\end{stegtabell}

Kontroll med $a = 2$: $\frac{(2\cdot 8)^2 \cdot 2}{4\cdot 32} = \frac{512}{128} = 4 = 2^2$. Det stemmer.
\end{eksempel}

\begin{vanligfeil}
\begin{itemize}[leftmargin=1.2em,itemsep=1pt,topsep=0pt]
  \item $a^3 \cdot a^4 = a^{12}$ gjelder ikke generelt. Med $a=2$ er $2^7=128$, men $2^{12}=4096$. Riktig: $a^7$.
  \item $(2a)^3 = 2a^3$ gjelder ikke generelt: $a=1$ gir $8\neq2$. Begge faktorene opphøyes: $(2a)^3 = 8a^3$.
  \item $-2^4 = 16$ er feil. Potensen regnes først: $-2^4 = -16$, men $(-2)^4 = 16$.
  \item $(a+b)^2 = a^2 + b^2$ gjelder ikke generelt: $a=b=1$ gir $4\neq2$. Reglene gjelder produkter, ikke summer.
\end{itemize}
\end{vanligfeil}

\nivaahode{1}

\oppg{1.1}{1}{\uh}
Skriv som én potens.
\begin{tasks}(4)
  \task $3^4 \cdot 3^2$
  \task $\dfrac{5^7}{5^3}$
  \task $(2^3)^4$
  \task $a^5 \cdot a$
\end{tasks}

\oppg{1.2}{1}{\uh}
Regn ut.
\begin{tasks}(4)
  \task $(2 \cdot 5)^3$
  \task $\left(\dfrac{2}{3}\right)^2$
  \task $(-5)^2$
  \task $-5^2$
\end{tasks}

\oppg{1.3}{1}{\uh}
Vi skal forenkle $\dfrac{a^4 \cdot a^3}{a^2}$.
\begin{tasks}(2)
  \task Skriv telleren som én potens.
  \task Forenkle hele brøken.
\end{tasks}

\nivaahode{2}

\oppg{1.4}{2}{\uh}
Forenkle.
\begin{tasks}(4)
  \task $(3x^2)^3$
  \task $\dfrac{6a^5b^2}{2a^2b}$
  \task $\dfrac{(2x^3)^2}{x^4}$
  \task $\left(\dfrac{x^2}{y}\right)^3 \cdot y^4$
\end{tasks}

\oppg{1.5}{2}{\uh}
Skriv om til samme grunntall. Regn deretter ut.
\begin{tasks}(3)
  \task $\dfrac{8^2}{2^5}$
  \task $\dfrac{9^3}{3^5}$
  \task $\dfrac{4^5}{8^3}$
\end{tasks}

\oppg{1.6}{2}{\uh}
Finn \textbf{den første feilen}. Løs riktig.
\[ \dfrac{(2b^3)^3\cdot b^2}{4b^6}\to\dfrac{8b^9\cdot b^2}{4b^6}\to\dfrac{8b^{18}}{4b^6}\to2b^{12} \]

\nivaahode{3}

\oppg{1.7}{3}{\uh}
Er påstanden alltid, noen ganger eller aldri sann for alle tillatte verdier av $a$ og $n$? Begrunn svaret.
\begin{tasks}(4)
  \task $a^2 \cdot a^3 = a^6$
  \task $(-a)^3 = -a^3$
  \task $\dfrac{a^6}{a^2} = a^3$
  \task $3^n + 3^n = 3^{n+1}$
\end{tasks}

\oppg{1.8}{3}{\uh}
Skriv $2^{12}$ som en potens med grunntall 4, 8 og 16.
Kan du skrive $2^{12}$ som en potens av 32 med heltallig eksponent? Forklar.

\begin{sjekk}
\begin{tasks}(2)
  \task Skriv som én potens: $\dfrac{x^8}{x^2} \cdot x$.
  \task Regn ut $(-3)^2$ og $-3^2$.
  \task! Hvorfor er $2^3 \cdot 2^4 = 2^7$? Forklar med faktorer.
\end{tasks}
\end{sjekk}

% ---------------------------------------------------------------------

\del{2}{Null, negative eksponenter og standardform}

\begin{ordliste}
negativ eksponent & En eksponent mindre enn null, for eksempel $5^{-2}$. & \mittsprak\\
tierpotens & En potens med grunntall 10, for eksempel $10^{-3}$. & \mittsprak\\
standardform & Et tall skrevet som $k \cdot 10^n$, der $1 \le k < 10$ og $n$ er et helt tall. & \mittsprak\\
\end{ordliste}

\begin{mustkunne}
\textbf{Kvotientregelen gir oss nye eksponenter.} Se på to måter å regne ut samme brøk:
\[
\frac{a^3}{a^3} = a^{3-3} = a^0 \quad\text{og}\quad \frac{a^3}{a^3} = 1
\qquad\qquad
\frac{a^2}{a^5} = a^{2-5} = a^{-3} \quad\text{og}\quad \frac{a\cdot a}{a\cdot a\cdot a\cdot a\cdot a} = \frac{1}{a^3}
\]
Derfor definerer vi:
\[ a^0 = 1 \qquad\qquad a^{-n} = \frac{1}{a^n} \qquad\qquad (a \neq 0) \]
\textbf{Standardform:} $k \cdot 10^n$ med $1 \le k < 10$. For eksempel er $45\,000 = 4,5 \cdot 10^4$ og $0,00032 = 3,2 \cdot 10^{-4}$.
\end{mustkunne}

\begin{eksempel}[1 — negative eksponenter]
Forenkle $\dfrac{2x^{-3} \cdot x^5}{x^0 \cdot 4x}$. Skriv svaret uten negative eksponenter.

\begin{stegtabell}
$= \dfrac{2x^{2}}{1 \cdot 4x}$ & Produkt: $x^{-3+5} = x^2$. Og $x^0 = 1$.\\
$= \dfrac{2}{4} \cdot x^{2-1} = \dfrac{x}{2}$ & Del tallene og potensene hver for seg.\\
\end{stegtabell}
\end{eksempel}

\begin{eksempel}[2 — regne i standardform]
Regn ut $(3 \cdot 10^5) \cdot (4 \cdot 10^{-8})$. Svar i standardform.

\begin{stegtabell}
$= (3 \cdot 4) \cdot 10^{5 + (-8)}$ & Gang tallene for seg og tierpotensene for seg.\\
$= 12 \cdot 10^{-3}$ & Dette er ikke standardform, fordi $12 > 10$.\\
$= 1,2 \cdot 10^{1} \cdot 10^{-3} = 1,2 \cdot 10^{-2}$ & Skriv $12 = 1,2 \cdot 10^1$.\\
\end{stegtabell}
\end{eksempel}

\begin{vanligfeil}
\begin{itemize}[leftmargin=1.2em,itemsep=1pt,topsep=0pt]
  \item $2^{-3} = -8$ er feil. En negativ eksponent gir en brøk: $2^{-3} = \frac{1}{8}$.
  \item $5^0 = 0$ er feil. $5^0 = 1$.
  \item $12 \cdot 10^{-3}$ er ikke standardform. Tallet foran må være minst 1 og mindre enn 10.
\end{itemize}
\end{vanligfeil}

\nivaahode{1}

\oppg{2.1}{1}{\uh}
Regn ut.
\begin{tasks}(4)
  \task $7^0$
  \task $6^{-2}$
  \task $10^{-2}$
  \task $\left(\dfrac{1}{2}\right)^{-1}$
\end{tasks}

\oppg{2.2}{1}{\uh}
Skriv i standardform.
\begin{tasks}(3)
  \task $52\,000$
  \task $0,0071$
  \task $340 \cdot 10^5$
\end{tasks}

\oppg{2.3}{1}{\uh}
Skriv uten negative eksponenter.
\begin{tasks}(3)
  \task $x^{-2}$
  \task $3a^{-1}$
  \task $\dfrac{1}{b^{-4}}$
\end{tasks}

\nivaahode{2}

\oppg{2.4}{2}{\uh}
Forenkle. Skriv svaret uten negative eksponenter.
\begin{tasks}(4)
  \task $\dfrac{a^{-2} \cdot a^6}{a^3}$
  \task $(2x^{-1})^3$
  \task $\dfrac{6x^2y^{-1}}{3x^{-1}y^2}$
  \task $\left(\dfrac{a^{-1}}{b^2}\right)^{-2}$
\end{tasks}

\oppg{2.5}{2}{\uh}
Regn ut. Svar i standardform.
\begin{tasks}(3)
  \task $(2 \cdot 10^3) \cdot (4,5 \cdot 10^{-7})$
  \task $\dfrac{6 \cdot 10^4}{1,5 \cdot 10^{-2}}$
  \task $(3 \cdot 10^{-2})^3$
\end{tasks}

\oppg{2.6}{2}{\mh}
Lyset bruker omtrent 8,3 minutter fra sola til jorda. Lysfarten er $3,0 \cdot 10^8$ m/s.
Finn avstanden fra sola til jorda i kilometer. Svar i standardform med to gjeldende siffer.

\nivaahode{3}

\oppg{2.7}{3}{\uh}
Er påstanden alltid, noen ganger eller aldri sann for alle tillatte verdier? Begrunn svaret.
\begin{tasks}(1)
  \task $2^{-n} < 0$ for heltall $n$
  \task $a^{-1} < a$
  \task Produktet av to tall i standardform er i standardform.
\end{tasks}

\oppg{2.8}{3}{\uh}
Hvilket tall er størst: $2^{-10}$, $10^{-3}$ eller $4^{-4}$? Begrunn uten kalkulator.

\begin{sjekk}
\begin{tasks}(2)
  \task Regn ut $3^{-2}$.
  \task Skriv $0,00045$ i standardform.
  \task! Er $25 \cdot 10^3$ skrevet i standardform? Forklar.
\end{tasks}
\end{sjekk}
